\chapter{Clifford algebras and spinors}\label{app:clifford}
This appendix provides the algebra behind the Dirac operator and the charge conjugation of \cref{ch:spectral_action,ch:almost_comm}.

\section{Clifford algebras}
The Clifford algebra $\mathrm{Cl}_n$ is generated by $\gamma^1,\dots,\gamma^n$ with
\begin{equation}
\{\gamma^\mu,\gamma^\nu\}=2\delta^{\mu\nu}.
\end{equation}
Its complexification is a matrix algebra:
\[
\mathrm{Cl}_{2m}\otimes\mathbb C\cong M_{2^m}(\mathbb C),\qquad
\mathrm{Cl}_{2m+1}\otimes\mathbb C\cong M_{2^m}(\mathbb C)\oplus M_{2^m}(\mathbb C).
\]
In even dimension there is a unique irreducible representation on \emph{spinors} $\mathbb C^{2^m}$, and the \emph{chirality} $\gamma_5\propto\gamma^1\cdots\gamma^{2m}$ anticommutes with every $\gamma^\mu$ and squares to $1$, giving a $\mathbb Z_2$ grading.

\begin{example}[explicit matrices]
For $n=2$: $\gamma^1=\sigma_1$, $\gamma^2=\sigma_2$, $\gamma_5=-\ii\gamma^1\gamma^2=\sigma_3$. For $n=4$: $\gamma^1=\sigma_1\otimes\id$, $\gamma^2=\sigma_2\otimes\id$, $\gamma^3=\sigma_3\otimes\sigma_1$, $\gamma^4=\sigma_3\otimes\sigma_2$, with $\gamma_5=\sigma_3\otimes\sigma_3$.
\end{example}

\section{Spin group and Dirac operator}
The bilinears $\tfrac14[\gamma^\mu,\gamma^\nu]$ span the Lie algebra of $\mathrm{Spin}(n)$, the double cover of $SO(n)$. On a spin manifold $M$ the spinor bundle $S$ carries the spin connection $\nabla$, and the \emph{Dirac operator}
\begin{equation}
D=\ii\gamma^\mu\nabla_\mu
\end{equation}
is self-adjoint on $L^2(M,S)$ with discrete spectrum and compact resolvent for compact $M$. In even dimension $D$ anticommutes with $\gamma_5$, so its spectrum is symmetric about $0$ with eigenspaces exchanged by chirality (compare the $\pm n$ pairs on $S^2$ in \cref{app:heat}). The Lichnerowicz formula $D^2=\nabla^*\nabla+\tfrac R4$ is what links the Dirac spectrum to scalar curvature.

\section{Charge conjugation and KO-dimension}
Spinors admit an antiunitary \emph{charge conjugation} $J$. Its properties depend only on the dimension modulo $8$, the \emph{KO-dimension}: with $J^2=\epsilon$, $JD=\epsilon'DJ$ and, in even dimension, $J\gamma_5=\epsilon''\gamma_5J$,
\begin{center}\small
\begin{tabular}{@{}l*{8}{c}@{}}
\toprule
$n \bmod 8$ & 0 & 1 & 2 & 3 & 4 & 5 & 6 & 7\\
\midrule
$\epsilon$ & $+$ & $+$ & $-$ & $-$ & $-$ & $-$ & $+$ & $+$\\
$\epsilon'$ & $+$ & $-$ & $+$ & $+$ & $+$ & $-$ & $+$ & $+$\\
$\epsilon''$ & $+$ & & $-$ & & $+$ & & $-$ & \\
\bottomrule
\end{tabular}
\end{center}
These are the same signs that define a real spectral triple in \cref{ch:almost_comm}. They also explain a Standard Model fact: the finite space has KO-dimension $6$, so the product with a four-dimensional manifold has total KO-dimension $4+6\equiv2\pmod 8$.

\mnote[Periodicity]{$\mathrm{Cl}_{n+8}\cong\mathrm{Cl}_n\otimes M_{16}(\mathbb R)$ (Bott periodicity), which is why the KO-dimension is defined modulo 8.}
